\chapter{Implementation contracts}
\label{ch:implementation}

\section{Reference implementation and component boundaries}

The initial implementation uses Python for orchestration and a deterministic
reference evaluator, exact decimal/rational arithmetic for booked amounts,
SQLite for transactional metadata and versioned records, and columnar files
for large scenario outputs. DuckDB is an optional analytical reader. These
are convenience choices justified by inspectability and modest initial scale;
they do not establish production performance. The IR and conformance tests
allow a later compiled executor without changing accounting meaning.

\begin{lstlisting}
accdsl/
  evidence/       documents, spans, filing contexts, extraction versions
  model/          typed IR, entities, books, symbols, contracts
  language/       grammar, parser, resolver, type checker, lowering
  constraints/    propagation, solver adapters, bounds and conflict cores
  accounting/     rule selection, posting validation, closing, consolidation
  economics/      contract and behavioural plugin interfaces
  execution/      scheduler, branch state, stochastic streams, run records
  reporting/      statements, notes, movement schedules, explanations
  storage/        versioned records, migrations, result readers
  interfaces/     Python API, CLI, structured import/export
  packages/       reviewed demonstration and accounting modules
\end{lstlisting}

The first service boundary is local and synchronous. A job abstraction can be
added for extraction and simulation without requiring distributed execution.
The model compiler returns an immutable execution plan; the executor receives
that plan and run inputs. Neither should retrieve live market data implicitly.
Input datasets and rule packages are selected and versioned before execution.

\section{Persistent records}

\begin{longtable}{L{33mm}L{103mm}}
\toprule Relation & Key information and constraints\\\midrule\endhead
\code{document} & Hash, original path/URI, filing ID, entity, publication and
availability times, media type, revision links. Original bytes are retained.\\
\code{source\_span} & Document hash, page/element locator, bounding box,
text/table representation, extraction run ID.\\
\code{observation} & Context ID, value or missingness tag, precision domain,
source-span IDs, semantic-reading ID, review status.\\
\code{context} & Entity/perimeter, concept, basis, currency/unit, instant or
interval, dimensions and taxonomy version. Contexts are immutable.\\
\code{assertion\_version} & Logical assertion ID, payload hash, economic
validity, ingestion time, public availability, superseded version and reason.\\
\code{model\_node} & IR node ID, type, serialised payload, dependency IDs,
package/model version; no executable arbitrary text.\\
\code{constraint} & Expression, scope, source/assumption references, supported
solver fragment, applicable alternative condition.\\
\code{alternative} & Parent model, added conditions, rationale, evidence,
coverage declaration and rejection history.\\
\code{economic\_event} & Stable event ID, contract/cohort, dates, quantities,
dependencies, scenario and branch.\\
\code{journal\_batch} & Event ID, book/basis, posting date, rule version,
balanced postings, measurement IDs and atomic commit status.\\
\code{run} & Evidence and model hashes, cutoff, scenarios, branch selection,
weight model, solver/numerical settings, seeds and environment.\\
\code{report\_cell} & Run, branch/scenario, context, amount or result family,
dependency references and completeness status.\\\bottomrule
\end{longtable}

Persistent identifiers are opaque; display labels can change without changing
identity. References use foreign keys or validated content-addressed links.
The full logical key of a reported fact includes its dimensions, not merely
company, label and year. Duplicate assertions are stored as distinct source
claims until a documented duplicate-resolution operation relates them.

Version records are append-only. A derived transaction-time interval can be
closed by the next version event without deleting the original payload.
Validity intervals can overlap when assertions conflict. Such overlap is
diagnosed and resolved by the evidence view, not prohibited by a database
constraint that would erase one assertion. Scenario states use structural
sharing of immutable historical nodes and branch-local changes.

Large simulation output is indexed by run, explanation, parameter draw,
scenario, innovation path, entity, book, date and account. Reducing this to
means across all dimensions would destroy the product's central feature.
Online summaries may compress draws within a declared distribution, while
retaining conditional summaries and representative replay identifiers. A
retention policy that removes individual traces must state which explanations
can still be reconstructed.

\begin{minipage}{\linewidth}
\section{Public API and result types}

The Python and command-line interfaces expose the same operations:
\begin{lstlisting}
ingest(manifest, extractor_config) -> EvidenceVersion
compile(source_or_ir, evidence_view, package_lock) -> ModelPlan | Diagnostics
reconcile(plan, cutoff, solver_config) -> HistoryFamily | Diagnostics
condition(family, assumptions) -> ConditionalFamily | Infeasible
simulate(family, scenario_set, run_config) -> RunResult
render(run_id, report_request) -> ReportBundle
explain(node_id, evidence_view_or_run_id) -> ExplanationGraph
compare(run_a, run_b, alignment_spec) -> Comparison
\end{lstlisting}
\end{minipage}

\code{HistoryFamily} contains a model-scope declaration, shared variables,
constraints, finite choices, feasibility status and solver-completeness status.
It may contain witness histories without claiming that those witnesses
enumerate the family. \code{RunResult} is a tagged union of concrete outputs,
conditional samples, certified bounds, explored ranges and unresolved outputs.
Probability-bearing results reference a probability-model record. Results
without that reference cannot be presented as probabilities.

Comparison requires an alignment specification: entity/perimeter, dates,
reporting basis, scenarios and shared innovations where applicable. Comparing
IFRS and US GAAP books holds economic events fixed and changes accounting
rules. Comparing management responses holds the specified scenario and
initial information fixed and changes the behavioural model. These are
different questions and must not be combined into an unexplained difference.

\section{Plugin interfaces}

Every plugin declares a version, supported types, input information dates,
determinism, numerical measure and capabilities. A contract plugin implements
\code{schedule}, \code{observe}, \code{transition} and \code{payoff}; it may
return future conditional events. A behaviour plugin implements
\code{decide} and declares which state and scenario information it uses.
A measurement plugin implements \code{measure} and returns an amount or
vector with a basis and numerical status. A reporting plugin maps classified
journals and balances to output contexts.

The plugin response must include all inputs that materially affect results.
An external calculator's global evaluation date, market-curve cache, calendar
version or random seed cannot remain invisible. For example, QuantLib's
floating coupon interface makes fixing-date conventions and pricing objects
relevant inputs; the adapter records them rather than relying on process-global
state. ACTUS risk-factor observations receive explicit dates and identifiers.

The constraint interface accepts a typed formula and a requested operation:
feasibility, witness, enumeration of declared finite choices, or bound. It
returns \code{sat}, \code{unsat} or \code{unknown}, plus evidence appropriate
to the operation. Timeouts and unsupported theories return \code{unknown}.
A witness proves existence within the formula, not completeness of a recovered
historical family.

\section{Human review and model-assisted extraction}

The review interface should put the source passage beside the proposed
observation, context and reconciliation impact. An analyst can accept a
mapping, add an alternative, reject an extraction, define a symbolic unknown
or request more evidence. These actions create structured changes with
reasons. The useful unit of review is a coherent disclosure and its links,
not an isolated number detached from its note.

A language model can identify candidate tables, draft concept mappings,
propose contract terms and suggest alternative readings. Its output must
reference source spans and a schema. Deterministic checks reject fabricated
source references, invalid units, inconsistent periods and unbalanced
postings. A model-generated interpretation that passes those checks remains
a proposal until its evidentiary status is established. Token probabilities
and self-reported confidence are never automatically used as weights over
economic histories.

For routine local research use, a lightweight interface and immutable source
versions are sufficient. Multi-user approval hierarchies, cryptographic
attestations and deployment security systems are outside this initial design.
They can be added for a concrete operational requirement without changing the
scientific representation.

\section{Failure recovery and reproducibility}

Extraction jobs and simulation jobs write explicit progress records and atomic
results. Restarting a failed document conversion reuses its document hash and
creates either the same idempotent extraction result or a new version under
changed settings. Restarting a simulation resumes only from a committed state
whose model and input hashes match the run. Numerical failures retain the
inputs needed for a focused diagnostic.

A run manifest records code revision when available, package lock, source
hashes, cutoff rule, model hash, scenario data, priors or weights, seed/stream
scheme, solver settings, decimal and rounding conventions, calendars,
environment and result locations. Absence of a usable Git revision is recorded
as unavailable with content hashes, rather than inventing a commit.

\section{Review decisions have consequences for later results}

A proposed reading starts in \code{Proposed}. Acceptance records a reviewer,
date, reason, verified locator and context, together with the reconciliation
result. Acceptance means that the stated layer has been reviewed: accepting
an extraction of a company's sentence does not accept its accounting
treatment. Competing defensible readings enter \code{AcceptedAlternative}
under a choice group. Rejection preserves the source and reason.
A correction creates a new reading; replacing an accepted reading requires
an explicit \code{Supersede} action. A prior reading therefore remains
reproducible without entering the latest model.

Suppose OCR reads Northstar's original 120 as 720. The source span supports
120; correcting extraction creates a new reading linked to the erroneous
one. The later revenue correction from 120 to 110 is different: it changes
the company's reported assertion. The original extraction remains faithful
to its document. Finally, a reviewer may accept that two notes conflict
without accepting either as a unique model constraint. The view then contains
a conflict set requiring alternatives or further evidence.

A changed reading or mapping invalidates descendant constraints, families,
measurements, reports and cached runs for a new view. Unaffected dependencies
can be reused. Old frozen runs retain their original decisions and evidence
IDs; they are marked stale for current use, not overwritten. A journal derived
under the old view is never silently attached to a new source. The local
workflow needs no multi-user approval hierarchy to enforce these rules.

\section{Identity must survive harmless reordering}

The reference serialisation uses UTF-8 JSON with NFC strings, sorted ASCII
keys, no insignificant whitespace and no floating-point literals. Exact
numbers are reduced rational numerator/denominator strings; zero is $0/1$.
Duplicate keys, NaN, infinities and unpaired Unicode surrogates are rejected.
Arrays preserve order unless the schema declares a set, such as a constraint
conjunction or source-ID set, in which case IDs are sorted. This normalisation
recognises representational equality, not arbitrary economic equivalence.

SHA-256 with a domain/version prefix identifies three different closures.
Evidence identity includes source bytes, locations, extraction and review
versions. Model identity includes typed IR, assumptions, rules, calculator
locks and resolution, while excluding formatting and source byte offsets.
Run identity adds the evidence view, horizon, scenarios, outputs, calendars,
numerical policy, random-stream scheme and environment. A corrected
source can change evidence and run identity without changing the numerical
model; a new rate assumption changes the model as well.

An economic event ID derives from owner, schedule version, date, phase,
operation and occurrence ordinal. The idempotency key pairs run ID and event
ID. Identical payloads return the committed result; different payloads under
the same key produce \code{EventIdentityConflict}. Independent events
may reorder only when read/write sets are disjoint or a package proves
commutation, including funding constraints. Alphabetical ordering cannot
settle an economic timing ambiguity.

Stochastic keys contain seed, coupling group, entity, contract/cohort,
economic date, shock kind and draw index. Branch enumeration order is absent.
Shared coupling groups may use the same uniform draw with different marginal
transforms. The reference fixture maps the first 64 bits of a SHA-256 key
digest, interpreted as unsigned big-endian integer $v$, to
$(v+1/2)/2^{64}$. This fixes example replay; it does not validate a production
Monte Carlo generator. Research runs select an assessed counter-based
generator and record transforms and numerical environment. Inner measurement
draws have different shock domains from outer realised defaults.

Schema migration preserves old and new payloads and a field/meaning mapping.
A representation-only migration must demonstrate equivalent supported
semantics; a changed recognition rule creates a new model. Restart resumes
only a committed state whose locks and identities match. SQLite owns source,
review and run metadata transactions; content-addressed files hold immutable
bytes and JSON, while DuckDB reads versioned analytical exports. Each mutable
transition has one owner.

\section{Adapters exchange typed records}

Exact method arguments and result fields are specified in
\code{spec/core-contract.md}; \path{spec/core-ir.schema.json} supplies closed
structural definitions. Every method returns a success record or a
diagnostic with code, source IDs, entity/period, dependency path and retry
status. The principal sequence is \code{ingest}, \code{review},
\code{compile}, \code{reconstruct}, \code{simulate}, \code{report}, with
\code{bound}, \code{explain} and \code{replay} over frozen identities.

Arelle supplies candidate facts with QName, entity, context, period,
dimensions, unit, accuracy, nil status and locator. Duplicates survive for
review. Its archived \code{ModelInstanceObject} interface exposes these
properties; validation success does not establish economic truth. Docling
supplies page/table proposals, locations and conversion status. Missing
headers and failed pages remain unresolved. Z3 receives named rational-linear
assertions and finite choices; witnesses are checked against exact inputs and
unknown results retain their reasons. The reference checks exercise Z3;
the other selected adapters remain interface designs grounded in inspected
official source snapshots.

ACTUS receives contract terms, calendars and dated risk factors and returns
schedules, transitions and payoffs. AccDSL applies accounting separately.
QuantLib receives an explicit evaluation date, instrument, curves, fixings,
calendar/day-count and measure; it returns dated measurements with numerical
status. No process-global valuation date may leak between branches.
Beancount exports concrete journals for later differential checks; it does
not own symbolic families. Import failures retain request ID and source hash
so retries are idempotent. Licence/version evidence remains in the software
ledger; integration acceptance requires pinned releases or source digests.
